\section{Hardware-in-the-Loop Development}
\label{sec:hil}
HIL simulation is the established bridge between simulation and
flight for UAV autopilots~\cite{jung2007hil,cai2009design}, and the
bench lineage of our group runs from the Pegasus real-time HIL
architecture~\cite{bittar2012pegasus} through X-Plane-based attitude
HIL~\cite{bittar2014hil} to the fixed-wing autopilot bench of
Santos~\cite{santos2018}. This section describes the present
incarnation: a lockstep serial bench in which the embedded controller
is the auto-generated production code.

\subsection{Controller Discretization}
\label{sec:discretization}
Both controllers are discretized at $T_s = 10$\,ms directly in Simulink
(automatic conversion of the compensators to their discrete-time forms;
the washout filter via Tustin transformation), producing single-rate
discrete models suitable for code generation. The discretized PID was
verified against the continuous design on the nonlinear model, with an RMS
deviation below $0.1\%$ of each command amplitude --- discretization is
\emph{not} a significant error source at this sample rate, as quantified
in Section~\ref{sec:sensitivity}.

\emph{A practical note on output saturation:} applying the automatic
conversion to the complete control loop with the output saturation blocks
included did not yield a working discrete model. The procedure adopted was
to discretize the control law alone and to re-insert the saturations
\emph{after} discretization. In the deployment reported here the
saturation limits (throttle in $[0,1]$, surfaces at $\pm 25^\circ$) are
enforced in the embedded C code, mirroring the host-side limits exactly.

\subsection{HIL Architecture}
\label{sec:hil-arch}
Fig.~\ref{fig:hil-loop} shows the architecture. The discrete controller is
translated to C by Simulink Embedded Coder (\texttt{ert.tlc}, ARM
Cortex-M, double precision) --- the automatic-code-generation workflow
of model-based aerospace software
practice~\cite{erkkinen2006automatic,erkkinen2009model} --- and
embedded in an STM32 microcontroller running FreeRTOS --- a class of
target for which HIL validation remains an active
topic~\cite{laabousse2025advanced}. The host PC runs \emph{only} the nonlinear plant, the
reference generation, and the serial interface. Host and target exchange
one frame per controller period over a 460\,800-baud serial link:
16 single-precision floats to the target (states, filtered references, and
mode switches) and 4 floats back (throttle, elevator, aileron, rudder),
delimited by a fixed header (\texttt{0xDE 0xAD}) and terminator.

The link operates in \emph{lockstep}, the mechanism by which
simulator and autopilot advance strictly by exchanged messages to make
runs deterministic and reproducible --- standard practice in
simulation stacks such as PX4~\cite{px4lockstep}: the plant blocks at
each 10\,ms simulation step until the controller reply arrives, so
simulated time advances only upon a completed exchange. Every HIL sample therefore
corresponds one-to-one to a SIL sample --- the comparison is free of
timing jitter by construction.

Two aspects of the campaign design are worth stating explicitly, as they
are what make the SIL and HIL results directly comparable:
\begin{itemize}
  \item \emph{Single source of references.} One script generates both the
        SIL and the HIL runs of every flight case, so the two campaigns
        cannot drift apart in their reference definitions.
  \item \emph{Mode selection over the link.} The three flight-mode
        switches travel in the frame rather than being compiled into the
        firmware, so a single binary serves all flight cases of the PID
        deployment and the target is never reprogrammed mid-campaign
        (the LQRy deployment deviates from this, as noted in
        Section~\ref{sec:campaigns}).
\end{itemize}

\begin{figure}[!t]
\centering
\includegraphics[width=\linewidth]{hil_loop.png}
\caption{HIL loop in lockstep: nonlinear plant on the host, generated
controller code on the STM32, binary serial frames at each 10\,ms step.}
\label{fig:hil-loop}
\end{figure}

\subsection{Link and Chain Validation}
Before any campaign, each element of the chain was validated
independently: (i) serial echo at 460\,800 baud with zero errors;
(ii) a frame-level test in which the embedded controller response to 200
known input frames matched the simulated discrete controller to within
$4\times10^{-8}$ (the single-precision quantization of the bus); and
(iii) the discretization check of Section~\ref{sec:discretization}.

\subsection{Determinism of the Lockstep Loop}
\label{sec:determinism}
A lockstep HIL loop built from double-precision arithmetic on both sides,
with single-precision quantization only on the wire, has no source of
run-to-run variability: there is no timing race, no dropped-sample
interpolation, and no floating-point reduction whose order can change.
The loop should therefore be \emph{bit-deterministic}.

We verified this property directly. The generated LQRy controller was
compiled and executed on a microcontroller model \emph{different} from the
one used in an earlier execution of the same code (STM32F303 versus
STM32F411, different core configurations and clock trees), driving the same
host model. Over a full 300-second flight case, the two executions
produced \emph{bit-identical} outputs on all recorded channels (maximum
difference identically zero, Fig.~\ref{fig:replica}); repeating additional
flight cases of the campaign reproduced the same result.

\begin{figure}[!t]
\centering
\includegraphics[width=\linewidth]{replica_caso1.png}
\caption{Bit-exact reproducibility: the same generated controller code
executed on two different microcontroller models over a 300-second flight
case. Maximum difference across all channels is identically zero.}
\label{fig:replica}
\end{figure}

This is a useful property for a validation bench, and not merely a
curiosity. Because the loop contributes no noise of its own, \emph{any}
deviation observed between two runs is attributable to a deliberate change
in the setup, which makes the bench a measuring instrument rather than a
source of scatter. It also means that a campaign can be reproduced exactly
by a third party, on different hardware, from the archived code and
reference definitions --- a stronger form of reproducibility than the
usual statistical agreement. The sensitivity study of the next subsection
relies on this: each experiment changes exactly one factor, and the
resulting difference is entirely due to that factor.

This determinism also settles a practical concern of embedded
deployment: block \emph{execution order} in a graphical model can, in
principle, alter the generated real-time code with respect to the
simulation. Here, with both controllers auto-generated from single-rate
discrete models, the cross-MCU bit-exactness rules out target-side
nondeterminism, and the quantization-floor SIL--HIL residual of the
synchronized PID deployment (Section~\ref{sec:campaigns}) bounds any
execution-order effect at the level of the bus quantization itself.

\subsection{Sensitivity of SIL--HIL Agreement}
\label{sec:sensitivity}
Deploying a control law on an embedded target involves a number of
implementation choices whose effect on SIL--HIL agreement is rarely
quantified. Exploiting the determinism established above, we injected each
factor individually into the validated PID bench and measured the
resulting discrepancy against the clean SIL run
(Table~\ref{tab:sensitivity}). All experiments use flight case 1; the
baseline row is the unmodified bench.

\begin{table}[!t]
\caption{Sensitivity of SIL--HIL agreement: each factor injected
individually into the validated bench (RMS vs.\ the clean SIL run, flight
case 1).}
\label{tab:sensitivity}
\centering
\begin{tabular}{lcc}
\toprule
Factor & $\theta$ [$^\circ$] & $h$ [m] \\
\midrule
Baseline (unmodified bench) & 0.030 & 0.09 \\
S1: start-up desynchronization (7\,s) & 0.014 & 0.07 \\
S2: airspeed reference mismatch (3\,m/s) & 0.447 & 258 \\
S3: controller discretization & 0.003 & 0.001 \\
S4: 1\% frame loss & 0.032 & 1.30 \\
\bottomrule
\end{tabular}
\end{table}

The ranking is informative, and in part counter-intuitive:

\subsubsection{Reference consistency dominates (S2)} An airspeed reference
differing by 3\,m/s between the SIL and HIL campaigns --- a plausible
outcome whenever the two campaigns are configured from separate scripts
--- produces a 258\,m RMS altitude discrepancy. This is not a fidelity
loss of the HIL chain at all: it is a comparison between two different
flight conditions. It is, by two orders of magnitude, the largest effect
measured, which motivates the single-source-of-references design of
Section~\ref{sec:hil-arch}.

\subsubsection{Start-up desynchronization leaves a persistent residual (S1)}
Embedded controllers commonly hold their integrators in reset until the
communication link engages, so the target starts from the trim command
while the simulated controller starts from its initial states and produces
an initial transient. The two trajectories then differ over the first
seconds and, in modes where an outer loop is open, that difference is
integrated and never recovered. Synchronizing the two start-ups --- as
done in the PID deployment --- removes this term; leaving them
unsynchronized leaves a residual that no amount of link quality can
remove.

\subsubsection{Frame loss is benign at low rates (S4)} Discarding 1\% of
the controller replies at random --- each discarded reply is replaced by
holding the previous command for that step, so the lockstep timing is
preserved --- degrades altitude agreement to 1.30\,m
but leaves attitude essentially unaffected, since the actuator rate limits
and the airframe dynamics filter isolated command glitches. This is
reassuring for links less reliable than the one used here.

\subsubsection{Discretization is negligible (S3)} The usual first suspect
--- the continuous-to-discrete conversion of the control law --- accounts
for $0.003^\circ$ in attitude and $10^{-3}$\,m in altitude at
$T_s=10$\,ms, three orders of magnitude below the reference-consistency
term.

\subsection{Embedded Campaign Results}
\label{sec:campaigns}
Eight flight cases (all combinations of the three mode switches, with
doublets on attitude/heading and steps on airspeed/altitude, 300\,s each)
were executed in SIL and HIL for both controllers, from a single script
per controller.

\subsubsection{PID} Table~\ref{tab:pid-silhil} reports the RMS SIL--HIL
discrepancy ($t>12$\,s). The residual reaches the quantization floor:
pitch attitude within $0.011^\circ$ in the worst case, altitude within
2.75\,m only where the altitude loop is open (integrated drift), and
essentially zero elsewhere. Fig.~\ref{fig:pid-c1} illustrates case~1: the
SIL and HIL traces are indistinguishable at plot scale.

\begin{table}[!t]
\caption{PID: RMS SIL--HIL discrepancy per flight case ($t>12$\,s).}
\label{tab:pid-silhil}
\centering
\begin{tabular}{lcccc}
\toprule
Case & $\theta$ [$^\circ$] & $\VT$ [m/s] & $h$ [m] & $\phi$ [$^\circ$] \\
\midrule
1 & 0.0028 & 0.0009 & 0.089 & $\sim$$10^{-6}$ \\
2 & 0.0007 & 0.0002 & 0.004 & $\sim$$10^{-6}$ \\
3 & 0.0035 & 0.0011 & 2.75  & $\sim$$10^{-6}$ \\
4 & 0.0007 & 0.0002 & 0.004 & $\sim$$10^{-7}$ \\
5 & 0.0109 & 0.0059 & 0.402 & $\sim$$10^{-6}$ \\
6--8 & $<10^{-3}$ & $<10^{-4}$ & $<10^{-4}$ & $<10^{-6}$ \\
\bottomrule
\end{tabular}
\end{table}

\begin{figure}[!t]
\centering
\includegraphics[width=\linewidth]{pid_sil_hil_caso1.png}
\caption{PID, flight case 1: SIL vs.\ HIL on the nonlinear model.}
\label{fig:pid-c1}
\end{figure}

\subsubsection{LQRy} The embedded LQRy is the earlier
\emph{fixed-gain} delivery --- a softer gain set than the
gain-scheduled design of Section~\ref{sec:lqry}, registered in the
repository~\cite{repo} --- so the bench results validate the
deployment pipeline for that build; re-deploying the gain-scheduled
law on the bench is future work
(Section~\ref{sec:conclusion}). The LQRy campaign was executed under the same
protocol, with per-case firmware builds and scripted programming of the
target. Its SIL--HIL residual is larger --- from $0.23^\circ$ to
$3.57^\circ$ in pitch depending on the flight case --- and is explained by
the start-up term quantified in S1: this implementation holds the
embedded integrators in reset until the link engages, so the embedded and
simulated runs begin from different states. The effect is deterministic
and confined to the transient; the doublet responses themselves track
closely.
% Figura HIL LQRy dedicada — descomentar apos regeneracao em ingles:
% \begin{figure}[!t]\centering
% \includegraphics[width=\linewidth]{fig_lqry_hil.png}
% \caption{LQRy, flight case 1: SIL vs.\ the original HIL campaign vs.\
% the corrected HIL run (SIL-sourced references, scripted flashing).}
% \label{fig:lqry-hil}\end{figure}

\subsection{Controller Comparison on the Bench}
\label{sec:fair}
With both controllers on the same bench, the same plant, the same
references, and the same metric window, the dominant free variable is
the control law --- subject to the harness asymmetries disclosed in
Section~\ref{sec:protocol}, which the equalization ablation showed not
to drive the conclusions. Fig.~\ref{fig:justo} shows flight case~1: on
the bench, the PID tracks airspeed and pitch more tightly, while the
LQRy regulates altitude better when the airspeed loop is active ---
recalling that the embedded LQRy is the earlier fixed-gain build, so
these bench observations characterize that deployment rather than the
gain-scheduled law of Section~\ref{sec:sil-comparison}.

\begin{figure}[!t]
\centering
\includegraphics[width=\linewidth]{justo_caso1.png}
\caption{Embedded comparison, flight case 1: PID and LQRy under identical
conditions on the same bench.}
\label{fig:justo}
\end{figure}
